\section{Threats to Validity}
\label{sec:threats}

\textbf{Internal validity.} The prototype includes parsers for
NTA models, property files, diagnostic traces, repair sites,
and admissibility artifacts. Bugs in parsing, trace reconstruction, SMT
encoding, candidate application, or admissibility checking may affect
the results. We mitigate this threat by validating every accepted
candidate using the external model checker and by recording repair
reports for each run.

\textbf{Construct validity.} The evaluation uses injected clock-bound
faults. Such faults are appropriate for evaluating clock-bound repair,
but they do not cover all real modeling errors. Real faults may involve
wrong resets, missing transitions, incorrect synchronizations, data
updates, or erroneous properties. The results should therefore be
interpreted as evidence for clock-bound repair, not for arbitrary NTA
repair.

\textbf{External validity.} Although the industry dataset covers several
practical domains, it is still a finite set of open or reconstructed
models. It may not represent all industrial NTA models.
Larger proprietary models, richer data variables, and more complex
property specifications may introduce additional challenges.

\textbf{Baseline fairness.} The baseline is an iterative TARTAR-style
repair workflow adapted to the multi-property setting by adding
full-property regression and admissibility checking. This makes it
stronger than a one-shot baseline, but it is still not identical to the
original setting for which TARTAR was designed. We therefore interpret
the comparison as a comparison against a strong single-property repair
strategy, not as a claim that TARTAR itself is unsuitable in general.

\textbf{Property selection.} Industry-case properties were selected from
model descriptions, papers, embedded queries, and intended timing
requirements. Some properties may be incomplete or may not capture all
domain-level requirements. We mitigate this by requiring all reference
models to satisfy the selected property set and by filtering out trivial
guard-mirror properties when constructing the benchmark.

\textbf{Repair-space limitation.} The current method focuses on
numerical clock-bound edits in guards and invariants
\cite{kolbl2019clock}. Some failed cases likely require repair actions
outside this space. Extending the method to reset repair,
synchronization repair, clock-reference repair, and structural repair is
left for future work.
